\section{Требования к турнирам для аккредитации}

Для аккредитации турнира в рейтинге RR, турнир должен соответствовать следующим требованиям:

\begin{itemize}
	\item Количество игроков: \textbf{16 и более}.
	\item Количество игр: \textbf{4 и более}.
	\item Турнир должен быть \textbf{открытым}. Определение открытого турнира дано в Положении о рейтинговой системе и аккредитации турниров.
	\item Играются \textbf{ханчаны}, а за каждым столом \textbf{четыре} игрока.
	\item Правила и схема рассадок на турнире не должны существенно отклоняться от рекомендаций, описанных в Регламенте RRC. Список допустимых вариаций правил приведен ниже.
	\item Заявка на аккредитацию должна быть подана в соответствии с Положением о рейтинговой системе и аккредитации турниров.
\end{itemize}

Допустимые \textbf{вариации правил} (варианты, соответствующие текущему регламенту, выделены жирным; предпочитайте их по возможности):
\begin{itemize}
	\item Акадоры: \textbf{есть} / нет;
	\item Стартовые очки: \textbf{30000} / 25000;
	\item Ума: предпочтительно \textbf{15/5} либо 30/10. Допускается любая другая симметричная ума.
	\item Нагаши манган: \textbf{нет} / есть;
	\item Банкротство (досрочное прекращение игры при уходе игрока в минус): \textbf{нет} / есть;
	\item Абортивные ничьи: \textbf{есть, кроме тройного рона} / есть, все / нет;
	\item Атамаханэ: \textbf{нет} / есть;
	\item Чомбо: \textbf{-20000 после умы} / обратный манган;
	\item Якитори: \textbf{нет} / есть;
	\item Ренхо: \textbf{манган} / отсутствует.
\end{itemize}

Необычные \textbf{вариации рассадок}, используемые на турнире, могут также послужить поводом для недопуска турнира в рейтинг RR. Мы рекомендуем придерживаться следующих вариантов:
\begin{itemize}
	\item Рассадка на турнире полностью предопределенная. Организаторы заранее готовят рассадку на весь турнир и по ней происходят игры. Предопределенная рассадка позволяет глобально и заранее минимизировать пересечения между игроками, но требует предварительной работы от организаторов по ее формированию согласно количеству игроков.
	\item Рассадка на турнире полностью автоматическая и состоит из следующих этапов последовательно:
	\begin{itemize}
		\item 1 или 2 случайных рассадки подряд. На турнирах среднего размера стоит ограничиться одной случайной рассадкой в самом начале; на больших турнирах допускается играть по случайной рассадке первые две игры;
		\item Швейцарская рассадка на все игры в середине турнира;
		\item 2 или 3 интервальные рассадки в конце турнира. Выбор интервала оставляется на усмотрение организаторов. Для мелких однодневных турниров стоит ограничиться одной интервальной рассадкой в последней игре.
	\end{itemize}
\end{itemize}
